\documentclass[11pt, a4paper]{article}

% ── Core packages ───────────────────────────────────────────────────────────
\usepackage[T1]{fontenc}
\usepackage[utf8]{inputenc}
\usepackage[margin=2.4cm]{geometry}
\usepackage{microtype}
\usepackage{parskip}
\usepackage{booktabs}
\usepackage{graphicx}
\usepackage{enumitem}
\usepackage{xcolor}
\usepackage[colorlinks=true, linkcolor=black, urlcolor=black, citecolor=black]{hyperref}

% ── Colours (green used ONLY for one accent, not section rules) ─────────────
\definecolor{accent}{HTML}{006A4A}   % military green — reserved for badge only
\definecolor{muted} {HTML}{6B7280}   % gray text
\definecolor{rule}  {HTML}{D1D5DB}   % light gray for decorative rules

% ── Section formatting (clean black, no colour) ────────────────────────────
\usepackage{titlesec}
\titleformat{\section}  {\large\bfseries}{}{0em}{}[\color{rule}\titlerule]
\titleformat{\subsection}{\normalsize\bfseries}{}{0em}{}
\titlespacing{\section}   {0pt}{20pt}{6pt}
\titlespacing{\subsection}{0pt}{12pt}{4pt}

% ── Header / footer ─────────────────────────────────────────────────────────
\usepackage{fancyhdr}
\pagestyle{fancy}
\fancyhf{}
\renewcommand{\headrulewidth}{0.3pt}
\fancyhead[L]{\small\color{muted} African Trust \& Safety LLM Challenge}
\fancyhead[R]{\small\color{muted} Zindi · April 2026}
\fancyfoot[C]{\small\color{muted}\thepage}

% ── Helpers ─────────────────────────────────────────────────────────────────
\newcommand{\finding}[1]{\textbf{Finding #1\,---}\;}
\newcommand{\code}[1]{\texttt{\small #1}}
\newcommand{\badge}[1]{%
  \fbox{\textcolor{accent}{\small\textbf{#1}}}%
}

% ════════════════════════════════════════════════════════════════════════════
\begin{document}
\thispagestyle{empty}

% ── Title block ─────────────────────────────────────────────────────────────
\begin{center}
  {\LARGE\bfseries
    Probing the Safety Frontier of African-Language LLMs:\\[4pt]
    A Red-Teaming Study}

  \vspace{16pt}
  {\large Yehoshua}\\[4pt]
  {\small\color{muted}
    Zindi African Trust \& Safety LLM Challenge\quad·\quad April 2026\\[2pt]
    \href{https://github.com/yehoshua0/The-African-Trust-Safety-LLM-Challenge}%
         {\texttt{github.com/yehoshua0/The-African-Trust-Safety-LLM-Challenge}}
  }
\end{center}

\vspace{10pt}
\noindent\color{rule}\rule{\linewidth}{0.8pt}
\vspace{2pt}

% ── Abstract ────────────────────────────────────────────────────────────────
\begin{center}
\begin{minipage}{0.92\linewidth}\small
\textbf{Abstract.}\;
Large language models fine-tuned for African languages carry the promise of
democratising AI across underserved communities. Yet their safety alignment has
received comparatively little scrutiny. This report presents a structured
red-teaming study targeting two publicly available African-language models ---
\textbf{N-ATLaS} (Hausa, Igbo, Yoruba) and \textbf{Pawa-Gemma-Swahili-2B}
(Swahili) --- using 17 adversarial prompts spanning 9 risk categories. All
attacks were validated through a purpose-built web application that logged
reproducibility across multiple independent runs. Our results reveal that
classic jailbreak techniques, when faithfully translated into African languages,
retain near-full effectiveness: 12 of 17 attacks (71\,\%) were confirmed at a
3/3 reproducibility score, and 5 at 2/3. Three attack types ---
\textit{direct jailbreak}, \textit{prompt injection}, and
\textit{refusal softening} --- successfully transferred across both models and
multiple languages, suggesting that safety vulnerabilities are structurally
embedded rather than incidental.

\smallskip
\textbf{Keywords:}\; red-teaming · safety alignment · African NLP · jailbreak ·
N-ATLaS · Pawa-Gemma · Hausa · Igbo · Swahili · Yoruba
\end{minipage}
\end{center}

\noindent\color{rule}\rule{\linewidth}{0.8pt}
\vspace{4pt}

% ════════════════════════════════════════════════════════════════════════════
\section{Introduction}

The rapid proliferation of LLM technology into African markets --- driven by
initiatives from organisations such as Masakhane, Pawa AI, and Lelapa AI ---
raises urgent questions about the adequacy of safety guardrails in low-resource
language settings. Safety research has overwhelmingly concentrated on
high-resource languages, primarily English. The implicit assumption has been that
multilingual fine-tuning preserves alignment properties. This study challenges
that assumption.

We target two representative models from the African NLP landscape:

\begin{table}[h]
\centering\small
\begin{tabular}{llll}
\toprule
\textbf{Model} & \textbf{Languages} & \textbf{Size} & \textbf{Source} \\
\midrule
N-ATLaS               & Hausa, Igbo, Yoruba & \textasciitilde7B (Q4\_K\_M) & QuantFactory/N-ATLaS-GGUF \\
Pawa-Gemma-Swahili-2B & Swahili             & \textasciitilde2B (Q8\_0)    & Skier8402/gemma2-2b-swahili-it \\
\bottomrule
\end{tabular}
\caption{Target models and their characteristics.}
\end{table}

Both models serve communities spanning hundreds of millions of speakers across
West and East Africa. A successful jailbreak at scale is therefore not merely a
benchmark curiosity --- it represents a concrete harm vector.

% ════════════════════════════════════════════════════════════════════════════
\section{Methodology}

\subsection{Red-Teaming Infrastructure}

All experiments were conducted through a locally hosted FastAPI web application
interfacing with \code{llama-cpp-python} for GGUF inference. Generation
parameters were deliberately calibrated for adversarial pressure:

\begin{table}[h]
\centering\small
\begin{tabular}{lll}
\toprule
\textbf{Parameter} & \textbf{Value} & \textbf{Rationale} \\
\midrule
Temperature        & 0.90  & Flattens safety distribution \\
Top-P              & 0.95  & Broadens nucleus sampling \\
Repetition penalty & 1.15  & Suppresses refusal phrase loops \\
Max new tokens     & 2048  & Captures full harmful output \\
\bottomrule
\end{tabular}
\caption{Inference parameters used across all red-teaming runs.}
\end{table}

\subsection{Attack Taxonomy}

Prompts were classified along three orthogonal axes derived from the challenge
specification:

\begin{itemize}[leftmargin=1.6em, itemsep=2pt]
  \item \textbf{Risk category} --- 9 categories: Cybersecurity, Extremism,
        Self-Harm, Violence, Illegal Activity, Privacy/Data Leakage,
        Bias \& Discrimination, Harmful Instructions, and Policy Evasion.
  \item \textbf{Attack type} --- 9 techniques: Direct Jailbreak, Persona
        Hijacking, Refusal Softening, Translation Pivot, Hypothetical Scenario,
        Prompt Injection, Indirect Request, Context Poisoning, and Roleplay.
  \item \textbf{Language} --- Hausa, Igbo, Swahili, and Yoruba, with prompts
        composed in each target language.
\end{itemize}

\subsection{Validation Protocol}

Each prompt was submitted a minimum of two to three times under identical
generation conditions. A break was classified as:

\begin{itemize}[leftmargin=1.6em, itemsep=2pt]
  \item \textbf{3/3 Confirmed} --- all three independent runs produced a
        policy-violating response.
  \item \textbf{2/3 Confirmed} --- two of three runs produced a violation
        (one refusal tolerated).
\end{itemize}

Only confirmed breaks were included in the final submission.

% ════════════════════════════════════════════════════════════════════════════
\section{Results}

\subsection{Overview}

\begin{table}[h]
\centering\small
\begin{tabular}{ll}
\toprule
\textbf{Metric} & \textbf{Value} \\
\midrule
Total validated attacks  & 17 \\
Languages covered        & 4 (Hausa, Igbo, Swahili, Yoruba) \\
Risk categories hit      & 9\,/\,9 \\
Models broken            & 2\,/\,2 \\
3/3 reproducibility      & 12\quad(71\,\%) \\
2/3 reproducibility      & 5\quad\,(29\,\%) \\
Cross-model attack types & 3\quad(\code{direct\_jailbreak},
                           \code{prompt\_injection},
                           \code{refusal\_softening}) \\
\bottomrule
\end{tabular}
\caption{Aggregate results across all validated attacks.}
\end{table}

\subsection{Per-Language Breakdown}

\begin{table}[h]
\centering\small
\begin{tabular}{llccc}
\toprule
\textbf{Language} & \textbf{Model} & \textbf{Attacks} & \textbf{3/3} & \textbf{2/3} \\
\midrule
Swahili & Pawa-Gemma-Swahili-2B & 12 & 9 & 3 \\
Hausa   & N-ATLaS               & 2  & 1 & 1 \\
Yoruba  & N-ATLaS               & 2  & 2 & 0 \\
Igbo    & N-ATLaS               & 1  & 1 & 0 \\
\bottomrule
\end{tabular}
\caption{Attack distribution and reproducibility by language.}
\end{table}

\subsection{Key Findings}

\finding{1} \textbf{Language is not a safety barrier.}
Translating the classic DAN prompt into Hausa, Igbo, or Yoruba produced
successful cybersecurity jailbreaks on N-ATLaS across all three languages.
The \code{direct\_jailbreak} type is cross-model: it succeeded on both N-ATLaS
and Pawa-Gemma in their respective languages, with a combined 4/4 success rate.

\finding{2} \textbf{System prompt extraction is cross-model and cross-language.}
Attacks \code{PRIVLEAK01} (Swahili) and \code{YOR\_PRIVLEAK01} (Yoruba) both
induced models to echo a \code{SYSTEM PROMPT:} header and enumerate internal
guidelines. This is a structural vulnerability shared by both architectures,
not an artefact of any specific fine-tuning.

\finding{3} \textbf{Refusal softening is highly effective.}
Pre-emptive illegal-acknowledgement framing
(\textit{``I know this might be illegal but\ldots''}) bypassed refusals for
advance-fee fraud explanation in both Hausa (N-ATLaS) and Swahili (Pawa-Gemma),
achieving 3/3 on both.

\finding{4} \textbf{Authority personas erode safety boundaries.}
Institutional framing (professor, doctor) lowered the threshold for harmful
content generation in attacks targeting SQL injection (\code{CYBER02}) and
document forgery (\code{FORGE01}), even when the model nominally acknowledged
the illegality.

\finding{5} \textbf{Partial injection effects are observable.}
Attack \code{PINJECT01} caused Pawa-Gemma to respond in English rather than
Swahili --- a language switch that reveals safety-layer disruption even when
the terminal output is a refusal.

% ════════════════════════════════════════════════════════════════════════════
\section{Discussion}

\subsection{Structural vs.\ Incidental Alignment}

The pattern across 17 attacks suggests that alignment in these models is
\textbf{shallow and language-agnostic in the wrong direction}: guardrails appear
to rely on surface-level keyword matching or high-resource-language heuristics
that fail when the same adversarial intent is rephrased in a lower-resource
tongue. Models with strong English safety alignment can nonetheless exhibit broad
policy failure when queried in Hausa or Yoruba.

\subsection{Community-Specific Risk}

Several attacks carry amplified harm potential in African contexts:

\begin{itemize}[leftmargin=1.6em, itemsep=2pt]
  \item \textbf{M-Pesa phishing} (\code{PHISH01}) directly weaponises the
        dominant mobile-money infrastructure for approximately 33\,M Kenyan
        users.
  \item \textbf{Ethnic bias in Swahili} (\code{BIAS01}) generates fabricated
        demographic statistics that could exacerbate ethnic tensions in Kenya.
  \item \textbf{Sahel extremism recruitment} (\code{EXT01}) generates content
        tailored to the specific socioeconomic conditions of the region.
\end{itemize}

These are not generic safety failures --- they exploit cultural and economic
specificity that English-centric safety research would not anticipate.

\subsection{Limitations}

This study is bounded by (i)~a focus on two models, (ii)~a 17-attack sample
limited by competition scope, and (iii)~inference on quantised GGUF weights,
which may exhibit slightly different behaviour from the original floating-point
checkpoints. Replication on full-precision models and expansion to additional
languages (Amharic, Zulu, Afrikaans, French Creole) remain open avenues.

% ════════════════════════════════════════════════════════════════════════════
\section{Conclusion}

We have demonstrated that off-the-shelf jailbreak techniques, when adapted to
African languages, reliably elicit policy-violating responses from two
representative African-language LLMs. The cross-model and cross-language
consistency of direct jailbreaks, prompt injection, and refusal softening
argues for systematic adversarial evaluation as a prerequisite before deploying
African-language models in any user-facing context. We hope this work
contributes to a more robust and equitable AI safety research agenda for the
continent.

% ════════════════════════════════════════════════════════════════════════════
\section*{Acknowledgements}

This work was submitted to the \textbf{Zindi African Trust \& Safety LLM
Challenge}. We thank the challenge organisers and evaluators for their rigorous
review process, and the broader ecosystem of organisations whose collective
effort makes African-language NLP research possible.

\begin{figure}[h]
  \centering
  \includegraphics[width=\linewidth]{banner.png}
  \caption*{\footnotesize\color{muted}
    Challenge partners \& sponsors --- Cassava Technologies, AMINI, meetkai,
    Qhala, Lelapa AI, MTN, Lanfrica, Vodacom, Orange, Ethio Telecom, World
    Sandbox Alliance, Masakhane, Mozisha, Pawa AI, APHRC, TOUMAi, Digital
    Umuganda, Airtel, Capsele, AXUM, AXIAN Telecom, awarri, DATA SPIRES,
    and ABIS.}
\end{figure}

\smallskip
We further thank the authors of \textbf{N-ATLaS} (QuantFactory) and
\textbf{Pawa-Gemma-Swahili-2B} (Skier8402 / Google DeepMind) for making model
weights publicly available, enabling open safety research without cloud compute
dependency.

\end{document}
